\clearpage
\fancyhead[L]{\small CONNAISSANCE DES TEMS.}
\hypertarget{native-650-lower}{}
\expandafter\def\csname HistoricalNoteText37\endcsname{Die vorliegende Anzeige wurde im August des vergangenen Jahres geschrieben.}
\begin{center}\small Göttingische gelehrte Anzeigen. Stück 6. Seiten 49..56. 10. Januar 1828.\\[3pt]\rule{.72\linewidth}{.3pt}\end{center}

\emph{Connaissance des tems, ou des mouvemens célestes, à l'usage des astronomes et des navigateurs, pour l'an \textup{1829}. Publiée par le bureau des longitudes.} Paris 1826. Bei Bachelier. 382 Seiten. Großoktav.

\begingroup\baselineskip=13pt\relax
Der Aufbau des astronomischen Kalenders ist noch derselbe wie in den früheren Jahrgängen. Das Verzeichnis der geografischen Längen und Breiten, das etwa 2000 Orte enthält, bedarf wohl einer gründlichen Prüfung oder Überarbeitung. Nach einer Bemerkung, die in den letzten Jahrgängen regelmäßig erscheint, hat das Längenbüro eine besondere Kommission damit beauftragt, und wir können ihre Ergebnisse in Kürze erwarten. --- Die \emph{Zusätze} füllen 163 Seiten und enthalten folgende Artikel. Tafel von \textsc{Puissant} zur Abkürzung der Berechnung der mithilfe des Polarsterns beobachteten Polhöhen und Azimute. Die astronomischen Beobachtungen, die an vielen Punkten des trigonometrischen Netzes für die neue Karte von Frankreich angestellt wurden, gaben Anlass zur Berechnung dieser nicht unpraktischen Hilfstafeln. Sie beziehen sich auf keine bestimmte Deklination und ließen sich daher auch auf Beobachtungen anderer Sterne außerhalb des Meridians anwenden, sofern diese nicht zu weit vom Pol entfernt sind. --- Bericht von \textsc{Legendre} und \textsc{Mathieu} an die Königliche Akademie der Wissenschaften über einen Aufsatz von Herrn \textsc{Puissant} zur Bestimmung der Gestalt der Erde aus geodätischen und astronomischen Messungen. Der Aufsatz bezieht sich zunächst auf die in Frankreich gemessenen Abschnitte zweier Parallelkreise und enthält dem Bericht \SourceLeaf{651} zufolge Berechnungsmethoden für die dabei auftretenden Aufgaben. Da der Aufsatz erst noch im Mémorial du dépôt de la guerre gedruckt werden soll und der Bericht nicht ins Einzelne geht, können wir den Wert der von Herrn \textsc{Puissant} aufgestellten Regeln noch nicht beurteilen. Nur bei der Ausgleichung der Winkel eines Dreieckssystems, das zwei gemessene Grundlinien miteinander verbindet und bei dem der berechnete Wert der zweiten Grundlinie etwas vom gemessenen abweicht, wird das Verfahren von Herrn \textsc{Puissant} angeführt. Es läuft darauf hinaus, die Korrekturen an den Winkeln jedes Dreiecks proportional zum Fehler der Summe der drei Winkel anzusetzen. Die Berichterstatter bemerken, dass dieses Prinzip bequem anzuwenden, seine Gültigkeit aber noch nicht bewiesen sei. Wir fügen hinzu, dass sich die \emph{richtige} Lösung dieser Aufgabe nach den Grundsätzen der Wahrscheinlichkeitsrechnung leicht finden lässt und sich dadurch die Unzulässigkeit jenes Prinzips erweist. --- Über die beiden großen Ungleichheiten in der Bewegung von Jupiter und Saturn, von \textsc{Laplace}. Dieser Aufsatz betrifft diejenigen Teile dieser Ungleichheiten, die vom Quadrat der störenden Kräfte abhängen. Nach einer Bemerkung von Herrn \textsc{Plana} ist es unrichtig, zwischen diesem zweiten Teil der großen Gleichung in der Saturnbewegung und dem zweiten Teil der großen Gleichung in der Bewegung des Saturn dasselbe Verhältnis anzunehmen, das zwischen den Hauptteilen (die von der ersten Potenz der störenden Masse abhängen) besteht. \textsc{Laplace} leitet hier die Beziehung zwischen diesen Teilen genauer her, doch das Ergebnis ist mit den von \textsc{Plana} gefundenen Zahlenwerten unvereinbar. Dieser Gegenstand scheint also noch weiterer Untersuchung zu bedürfen, und es wäre wünschenswert, die gesamte Theorie der gegenseitigen Störungen von Jupiter und Saturn von Grund auf ganz anders als bisher zu behandeln. --- Bemerkungen desselben Verfassers über verschiedene Gegenstände der Mécanique Céleste betreffen (zum Teil als Erwiderungen auf einige Einwände von Herrn \textsc{Plana}) die Bewegung der Bahn des äußersten Saturnsatelliten, eine periodische Gleichung des Merkur und die Wirkung der Fixsterne auf das Sonnensystem. --- Über die Messung eines Bogens des Parallelkreises auf $45^\circ$ Breite, von \textsc{Brousseaud} und \textsc{Nicollet}. Dieser sehr interessante Aufsatz berichtet über die Operation, mit der in den Jahren 1822 und 1823 der astronomische Längenunterschied zwischen dem Mont-Cenis und dem Turm von Marennes an der Mündung der Garonne mithilfe von Pulversignalen bestimmt wurde. Der ganze Bogen umfasst 8 Längengrade, und die Pulversignale wurden an fünf Zwischenpunkten gegeben. Im Jahr 1822 wurde auf diese Weise der Längenunterschied zwischen dem Mont-Cenis und dem Puy d'Isson im Departement Puy de Dome durch das Zusammenwirken der erwähnten französischen Beobachter und der italienischen Astronomen \textsc{Carlini} und \textsc{Plana} bestimmt. Im Jahr 1823 bestimmten die französischen Beobachter allein den westlichen Abschnitt, der ganz auf französischem Gebiet liegt. Der Aufsatz enthält sowohl alle von französischer Seite angestellten astronomischen Beobachtungen als auch die daraus folgenden Ergebnisse für die Längenunterschiede. Die Ergebnisse wurden mithilfe der ähnlichen, in den Jahren 1821 und 1824 von den italienischen Astronomen ausgeführten und zum Teil bereits veröffentlichten Operationen bis zur Sternwarte von Padua ausgedehnt. So umfasst dieses große Unternehmen bereits einen Bogen von 13 Längengraden. Später wurde der Bogen noch bis Fiume erweitert, und es besteht Hoffnung auf eine weitere Verlängerung. Die geodätische Bestimmung des Bogens von Marennes bis Fiume beruht auf 106 Dreiecken erster Ordnung, von denen 90 von französischen, die übrigen von österreichischen und sardinischen Ingenieuren gemessen wurden. Es wäre sehr wünschenswert, dass auch diese Dreiecke bald ausführlich veröffentlicht werden. Wir haben dazu begründete Hoffnung, da nach einer Mitteilung am Schluss des Aufsatzes \emph{alles}, was sich auf die Messung dieses Parallelkreisbogens bezieht, im 9. Band des Mémorial du dépôt de la guerre gedruckt werden soll. Zu den im vorliegenden Aufsatz \SourceLeaf{652} enthaltenen Beobachtungen bieten sich noch einige Bemerkungen an. Die Zeitbestimmung stützte sich auf absolute Zenitdistanzen, die mit \textsc{Borda}s Repetitionskreisen an Sternen beobachtet wurden. Deren scheinbare Positionen wurden den Hilfstafeln \textsc{Schumacher}s entnommen. Da bei dieser Arbeit die größtmögliche Genauigkeit der Zeitbestimmung von höchster Bedeutung ist, muss man bedauern, dass an den Beobachtungsplätzen anstelle jenes Verfahrens keine transportablen Mittagsfernrohre aufgestellt wurden. Bei so wichtigen und weitreichenden Operationen hätten die dadurch entstehenden Mehrkosten keine Rolle spielen können. Bei den Zeitbestimmungen anhand von Sternen auf verschiedenen Seiten des Meridians zeigen sich Unterschiede von mehr als drei Zeitsekunden. Dabei bleibt es bedenklich, lediglich den Mittelwert zu nehmen. Auch scheint an einigen Punkten die Zahl der Abende, an denen die Pulversignale beobachtet wurden, zu gering, um völlige Gewissheit über die absolute Zuverlässigkeit der Ergebnisse zu geben. So wurde der Längenunterschied zwischen Solignat und La Jonchere an drei Abenden jeweils durch 10 Pulversignale bestimmt, nämlich

\begin{center}\small\setlength{\tabcolsep}{3pt}
\begin{tabular}{lrll}
August & 12 & . . . . . . . . . . . . . . . . . & $6^{\mathrm m}\ 47^{\mathrm s}86$ \\
 & 19 & . . . . . . . . . . . . . . . . . & $6\quad49.81$ \\
 & $22'$ & . . . . . . . . . . . . . . . . . & $6\quad50.17$ \\
\end{tabular}\end{center}
\noindent Man konnte durchaus keine Ursache ausfindig machen, die die beträchtliche Abweichung des ersten Ergebnisses von den beiden anderen erklärt hätte. Man schloss jenes ganz aus und hielt sich nur an das Mittel dieser beiden. Die Messung des Bogens zwischen La Jonchere und Saint Preuil beruht sogar nur auf den Beobachtungen der Pulversignale einer einzigen Nacht, da der Berg Puy Cogneux, auf dem die Signale gegeben wurden, in St. Preuil, wie es scheint, nur bei stärkerer terrestrischer Refraktion sichtbar war. Man war hier bei der Auswahl der Plätze nicht vorsichtig genug gewesen. Da die große Veränderlichkeit der terrestrischen Refraktion eine bekannte und bei großen geodätischen Operationen täglich ins Auge fallende Erscheinung ist, hätte man, bevor St. Preuil als Beobachtungsplatz gewählt wurde, allerdings die Möglichkeit der Gefahr berücksichtigen sollen, dass das dazwischenliegende hohe Gelände den Signalplatz bei geringerer oder mittlerer Refraktion verdecken könnte. Als Entschuldigung dienen allerdings das für so ausgedehnte Operationen zu kleine Personal und die beschränkte Zeit. Dennoch bleibt es auffallend, dass der in St. Preuil eingesetzte Beobachter anfangs sogar keine \emph{genaue} Kenntnis von der Richtung des Signalplatzes hatte, obwohl die Mittel dazu leicht zu finden sein mussten, da Puy Cogneux ein Hauptdreieckspunkt gewesen war. Die Länge eines Längengrades auf dem Breitenkreis bei $45^\circ\,43'\,12''$ Breite wird schließlich aus dem gesamten Bogen von Marennes bis Padua zu 77865,75 Meter berechnet. Aus dem Vergleich der sechs einzelnen Bestandteile ermittelt der Rezensent für dieses Ergebnis einen zu befürchtenden mittleren Fehler von 25 Metern. Angesichts der vorhin erwähnten Umstände dürfte die verbleibende Unzuverlässigkeit jedoch eher noch größer sein, und daher dürften die Ergebnisse, die für die Abplattung des Erdsphäroids aus dem Vergleich dieses Längengrades mit verschiedenen gemessenen Breitengraden gezogen werden, noch immer kein sehr großes Gewicht haben. — Über die Länge von Port-Praslin in Neu-Irland, von Fregattenkapitän Duperrey. Eine neue Berechnung der von Veron im Jahr 1768 beobachteten Sonnenfinsternis nach den Tafeln von \textsc{Burckhardt} und \textsc{Damoiseau} und eine vom Verfasser an Bord der Coquille vorgenommene chronometrische Bestimmung ergaben sehr gut übereinstimmende Ergebnisse. — Über ein Mittel, die Wirkung der Kapillarität bei Barometern aufzuheben, von \textsc{Laplace}. Da das Quecksilber im Gefäß zum Rand hin konvex ist, kann man den Stift, der die Oberfläche dieses Quecksilbers berührt und dessen unteres Ende dem Nullpunkt der Skala entspricht, so platzieren, dass die Wirkung der Kapillarität aufgehoben wird. Herr \textsc{Brunard} hat dazu eine Hilfstafel angegeben. — Über die Veränderung, welche die Konstante der \SourceLeaf{653} Präzession infolge des Königsberger Katalogs der Fundamentalsterne erfährt, von \textsc{Bessel}. Der Inhalt dieses Aufsatzes ist bereits aus den Astronomischen Nachrichten bekannt. --- Beobachtungen und Elemente des am 19. Mai 1825 entdeckten Kometen, von \textsc{Gambart}. --- Über die Anziehung der Sphäroide von \textsc{Poisson}. Dieser gehaltreiche Aufsatz besteht aus zwei Abteilungen. Die erstere beschäftigt sich mit der Entwicklung der Kugelfunktionen (so möchten wir die Funktionen zweier veränderlicher Größen, die allgemein jeden Punkt einer Kugelfläche bestimmen, nennen) in Reihen, die nach einem bekannten, von den Analysten vielfach behandelten Gesetz fortschreiten, und ist mit der diesem großen Geometer eigenen Eleganz durchgeführt. Der für das Fundamentaltheorem dieser Umwandlung bereits früher gegebene Beweis hat hier einen Zusatz erhalten, um einigen dagegen vorgebrachten Entwürfen zu begegnen. Dieser Zusatz behebt, richtig verstanden, allerdings alle Schwierigkeiten, obwohl eine vollständigere Unterscheidung der dabei vorkommenden unendlich kleinen Größen die Klarheit des Beweises noch vollkommener machen würde. Wenigstens zeigt das, was kürzlich\HistoricalNote{37} von einem berühmten englischen Geometer im Maiheft des Philosophical Magazine gegen die neue Darstellung von \textsc{Poisson} vorgebracht wurde, auch wenn es deren Wesen gar nicht trifft, dass diese Darstellung noch missverstanden werden konnte. Die zweite Abteilung enthält allgemeine Untersuchungen über die Anziehung beliebiger Körper, die mit großer analytischer Kunst ausgeführt sind und tiefer gehen als bisherige Untersuchungen. Doch beruht alles auf der Entwicklung in Reihen, und Ergebnisse erhält man nur insofern, als die Abweichung des Körpers von der Kugelgestalt als klein vorausgesetzt wird. Für den Fall, dass diese Bedingung wegfällt, ist nach den Worten des Verfassers ‘die Gestalt des Gleichgewichts von den Geometern noch nicht bestimmt, und wenn dies Ellipsoid die einzige Figur wäre, bei der das Gleichgewicht möglich ist, so würde die sonderbare Folgerung notwendig sein, dass das Gleichgewicht bei einer Geschwindigkeit der Rotationsbewegung unmöglich wird, die doch noch nicht diejenige ist, bei der die Flüssigkeit anfängt, sich zu zerstreuen.’ Wir würden lieber sagen, dass wir noch keinen mathematisch strengen Beweis dafür besitzen, dass das Ellipsoid die einzig mögliche Figur für das Gleichgewicht ist. Denn in der Tat enthält jene Folgerung durchaus nichts Widersprüchliches.
\par\clearpage\endgroup
